\section{Six Distinct Case Studies}\label{app:cases}
The cases below use the selection rule in Section~\ref{sec:test-quality}. ``Success'' denotes the official instance-level result. Raw before-fix and after-fix observations are reported separately when the log parser assigns results to different test names. Source paths and complete transition sets are indexed in \code{evidence.json} beside the paper.

We assess TRACE's generated tests using four properties: \textbf{P1, meaningful assertions}: the test explicitly checks an expected outcome; \textbf{P2, issue-related failure}: the test fails on buggy code because of the reported problem, not a typo, import error, or setup failure; \textbf{P3, passes after the fix}: the same test executes and passes after the reference fix; and \textbf{P4, requirement alignment}: the test uses a setup and assertions that represent the behavior requested in the issue. These are retrospective assessments from submitted patches and saved execution evidence, not outputs of an implemented automatic quality checker. P1--P4 are scored only for TRACE to assess the proposed workflow; baseline behavior is described for comparison but was not separately scored with this rubric. A passing post-fix test alone does not establish issue reproduction. ``Partial'' indicates that only part of the reported scenario is represented; ``not established'' indicates insufficient evidence, rather than a confirmed negative result.

\subsection{Shared Success: Astropy 14508}
\textbf{Issue.} FITS (Flexible Image Transport System) is a file format used in astronomy. Its header stores metadata in 80-character records called cards, containing a keyword, a value, and an optional comment. The HIERARCH convention allows keyword names longer than the usual eight characters, while retaining the 80-character record limit.\footnote{Astropy FITS header documentation: \url{https://docs.astropy.org/en/stable/io/fits/usage/headers.html}.}

In this issue, a decimal value is written with unnecessarily many digits, leaving too little space for the comment and causing it to be cut off. \textbf{Generated checks.} Both methods construct the reported card using the value 0.009125 and compare its serialized representation with the compact expected string. Baseline additionally checks the stored value and comment. \textbf{Before.} Both new tests fail with a VerifyWarning about comment truncation. \textbf{After.} Each passes, and each official report records one added fail-to-pass test without added persistent or pass-to-fail failures. \textbf{Relevance.} The tests check that the number is written compactly and the full comment is preserved, directly matching the problem described in the bug report. \textbf{Implication.} Both methods generate a native test that checks the expected behavior; TRACE's accepted issue specification does not produce an observable success advantage.
\begin{lstlisting}[language=Python,caption={Shared Astropy oracle, condensed from the submitted tests.}]
c = fits.Card("HIERARCH ESO IFM CL RADIUS", 0.009125,
              "[m] radius arround actuator to avoid")
assert str(c) == _pad(
    "HIERARCH ESO IFM CL RADIUS = 0.009125 / "
    "[m] radius arround actuator to avoid")
\end{lstlisting}
The comment spelling \code{arround} is reproduced verbatim from the issue example.

\textbf{TRACE quality assessment.} P1 met: the test explicitly compares the serialized card with the compact expected representation. P2 met: the buggy run encounters the reported comment-truncation warning while constructing the card, before reaching the final comparison; this is an issue-related failure, not a setup error. P3 met: the test passes after the fix. P4 met: the chosen value and comment directly exercise the reported formatting problem.

\subsection{Shared Success: Django 11163}
\textbf{Issue.} An explicitly empty fields list should produce an empty dictionary. \textbf{Generated checks.} Both workflows add the same test method to an existing test class that uses the database.
\begin{lstlisting}[language=Python,caption={Identical submitted Django 11163 addition.}]
def test_empty_fields(self):
    item = ColourfulItem.objects.create()
    self.assertEqual(model_to_dict(item, fields=[]), {})
\end{lstlisting}
\textbf{Before.} On buggy code, \code{model_to_dict(item, fields=[])} returns a dictionary containing model fields instead of an empty dictionary, so the assertion fails. \textbf{After.} The added test passes after the reference fix, and the official evaluator marks both workflows successful. It lists two fail-to-pass test entries, but only one is the newly generated test; the other is an existing many-to-many test. The reported count therefore does not mean that either agent wrote two tests. \textbf{Relevance.} The assertion checks exactly the behavior requested in the issue and uses a model already available in Django's test suite. \textbf{Implication.} Both workflows produced the same test patch, so TRACE's additional issue-specification step showed no benefit for this clear report in the observed runs. Nevertheless, the harness reported coverage deltas of 0.5 for baseline and 1.0 for TRACE. Because the patches are identical, patch differences cannot explain these coverage results.

\textbf{TRACE quality assessment.} P1 met: \code{assertEqual} checks for an empty dictionary. P2 met: the assertion fails because buggy code returns model fields for an explicitly empty field list. P3 met: the added test passes after the fix, as shown in the raw execution evidence. P4 met: the input and expected result directly match the issue. Parser attribution to an additional existing test does not change this assessment of the new test.

\subsection{Shared Failure: Django 12155}
\textbf{Issue.} When a docstring starts with text rather than a blank line, Django leaves extra indentation on later lines. This can cause errors when the documentation is displayed. \textbf{Generated checks.} Both workflows add assertions checking that indentation is removed correctly and that the docstring is split into the expected parts. TRACE also checks that the displayed documentation does not contain a formatting-error message from docutils, the documentation-processing library. \textbf{Before.} The saved execution logs show two new assertion failures for each workflow. These failures expose incorrect indentation handling; TRACE also detects a documentation-formatting error. \textbf{After.} The reference fix removes the \code{trim_docstring} helper, but the existing test module still imports it. That import fails before the runner can execute the newly added tests. The evaluator therefore cannot match those tests across the before-fix and after-fix runs, records no added fail-to-pass tests, and marks both workflows unsuccessful. \textbf{Relevance.} The assertions target the reported problem, but the tests rely on a helper that is no longer available after the fix. \textbf{Implication.} Exposing the bug before the fix is only part of a successful reproduction: the same test must also run and pass on the fixed code. Neither workflow demonstrated that here. Validating the issue specification did not prevent this test-compatibility problem, and the reference fix was not available to the agents during generation.

\textbf{TRACE quality assessment.} P1 met: the generated tests check whitespace handling and reject rendering-error output. P2 met for the observed new failures: retained indentation and rendering diagnostics concern the reported problem. P3 not met in this evaluation: the removed helper prevents the module from loading, so no passing post-fix execution of the new tests is demonstrated. P4 met for the intended indentation and rendering requirements, but reliance on the helper makes the submitted tests incompatible with the fixed repository.

\subsection{Shared Failure: Django 13513}
\textbf{Issue.} The report states that Django's debugging error page displays an earlier exception even when \code{raise ... from None} explicitly requests that it be hidden. This concerns Python's exception-context suppression. \textbf{Generated checks.} Both workflows create one exception while handling another and use \code{raise ... from None} to request that the earlier exception not be displayed. They then check the exception information collected for the debug report. \textbf{Before and after.} The new tests pass on both buggy and fixed code because Django already handles the particular scenario they check. Neither workflow produces a new test that fails before the fix and passes afterward, so both receive an unsuccessful result. A separate failure in an existing debug-template test is not a failure of these newly generated checks. \textbf{Relevance.} The tests address exception hiding, which is related to the report, but miss the specific case addressed by the official fix: an exception without a traceback, meaning that it carries no record of where it was raised. This detail was identified in our post-hoc review of the evaluation results and was not supplied to the agents for generation or repair. \textbf{Implication.} A test can contain sensible assertions about the right general topic yet fail to exercise the actual bug. TRACE's accepted issue specification did not prevent that outcome. Both methods receive a coverage delta of 0.647059 for this issue despite failing to reproduce it; executing relevant code is not the same as checking the broken behavior. The baseline agent also inspected repository history, which requires review against the experiment's restrictions.

\textbf{TRACE quality assessment.} P1 met: explicit assertions check the exception information. P2 not met: the generated tests pass on buggy code instead of exposing the bug. P3 met as an execution observation: the tests pass after the fix; because they also pass before it, this does not show successful reproduction. P4 partial: the tests check related exception-hiding behavior but omit the exception-without-a-traceback scenario. They therefore cannot distinguish the buggy version from the fixed version.

\subsection{Disagreement: Django 13344}
\textbf{Issue.} Middleware is code that processes a web request or response around a Django view. In the reported bug, the first middleware's \code{process_response()} method receives a coroutine, an asynchronous operation that still needs to be awaited, instead of a completed \code{HttpResponse}.

\textbf{Generated checks.} Both workflows add a helper middleware that records the object passed to \code{process_response()}. They then assert that the recorded object is an \code{HttpResponse}. The baseline places its helper before Django's existing \code{SecurityMiddleware}, following the report's example. TRACE instead places its helper before a new async-only middleware written by the agent. TRACE chose that setup after temporary experiments showed that it produced the reported coroutine symptom. Its issue specification had been accepted on the second extraction attempt, so this was not a fallback run.
\begin{lstlisting}[language=Python,caption={Expected response type checked by both workflows, simplified.}]
self.assertIsInstance(response, HttpResponse)
\end{lstlisting}

\textbf{Before.} Both new tests fail because their recording helpers receive a coroutine rather than an HttpResponse. \textbf{After.} The baseline test passes after the reference fix. TRACE's test still receives a coroutine and fails. The official evaluator therefore marks the baseline successful and TRACE unsuccessful.

\textbf{Reading the evaluation records.} The baseline's raw logs show the new test failing before the fix and passing afterward. However, the evaluator's log parser labels an existing test, \code{test_get_asgi_application}, as the fail-to-pass entry and cannot match the new test name between runs. We report the official outcome while noting that its test-name attribution differs from the raw logs.

\textbf{Relevance.} Both tests check the expected response type, but they use different middleware setups. TRACE reproduced a similar symptom without demonstrating that its setup exercised the same underlying bug. Its continued failure after the fix does not by itself identify the precise cause. \textbf{Implication.} Checking that the issue specification is supported by the report does not guarantee that the agent will build a faithful executable test. The response assertion and the surrounding setup both matter.

\textbf{TRACE quality assessment.} P1 met: the test explicitly checks for an HttpResponse. P2 not established: it fails with the reported symptom, but the evidence does not establish that the same underlying bug caused it. P3 not met: the test still fails after the reference fix. P4 partial: the expected response type matches the report, but the custom middleware setup differs from its example and is not shown to reproduce the same scenario.

\subsection{Disagreement: Django 15098}
\textbf{Issue.} Django does not correctly recognize language codes that specify both a writing system and a region. For example, \code{en-latn-us} means English, written using the Latin alphabet, for the United States. A valid web address using this language prefix can return a 404 Not Found response instead of the requested page. \textbf{Generated checks.} Both workflows add tests checking that the language-prefixed address works and that Django recognizes the language code. The baseline adds a language-lookup assertion to an existing test. TRACE creates a separate lookup test that also checks codes with mixed uppercase and lowercase letters. Before generation, TRACE tries three times to extract an issue specification that passes validation. None is accepted, so the agent receives only the original report, just as in the baseline. This run still counts as a TRACE result. Its recorded generation time includes the unsuccessful preprocessing attempts, and their API charges contribute to workflow spending; the agent-only cost figures exclude those requests.
\textbf{Before.} The URL tests and language-lookup assertions generated by both workflows fail. \textbf{After.} Both URL checks pass. The existing lookup test modified by baseline still fails, while TRACE's new lookup test passes in the raw log. Both suites retain additional existing failures or errors. The official baseline report records two added fail-to-pass entries but also an additional persistent failure relative to base, yielding an unsuccessful result. TRACE records three added fail-to-pass entries without an added persistent failure and succeeds. The parser cannot match TRACE's new lookup test name between runs, even though the raw after-fix log records it as passing.
\textbf{Relevance.} The URL tests directly check the reported behavior, while the additional language-lookup assertions check different aspects of language-code handling. \textbf{Implication.} The harness checks whether the generated patch introduces additional failures relative to its base comparison, not simply whether a new test changes from failing to passing. This TRACE-only success occurred without an accepted issue specification and therefore cannot demonstrate that an accepted issue specification caused the gain. Baseline's history inspection remains a separate limitation.

\textbf{TRACE quality assessment.} P1 met: explicit checks cover language-prefixed URLs and language lookup. P2 met for the observed generated checks: buggy execution fails on the language-code behavior rather than a syntax or setup error. P3 met in the inspected raw logs: the URL checks and new lookup method pass after the fix; parsed attribution of the lookup method remains uncertain. P4 met for the reported language/script/region behavior. TRACE receives official success, but because preprocessing fell back to the original report, these properties cannot be credited to an accepted issue specification.

\onecolumn
\section{Full Paired Outcomes}\label{app:outcomes}
Table~\ref{tab:all} retains the fixed assignment order used in the saved instance list, rather than sorting by repository or result. Success and Unsuccessful denote official instance-level outcomes.
\begin{table*}[t]
\caption{All Assigned Instances: Official Outcomes and Coverage Delta}
\label{tab:all}
\centering\scriptsize
\begin{tabular}{lccrr}
\toprule
Instance & Baseline & TRACE & Delta baseline & Delta TRACE\\
\midrule
\code{pydata__xarray-3677} & Success & Success & 1.000000 & 1.000000\\
\code{django__django-14089} & Success & Success & 0.500000 & 0.500000\\
\code{django__django-11163} & Success & Success & 0.500000 & 1.000000\\
\code{sympy__sympy-24539} & Success & Success & 0.285714 & 0.000000\\
\code{scikit-learn__scikit-learn-14894} & Success & Success & 0.416667 & 0.416667\\
\code{sphinx-doc__sphinx-9711} & Success & Success & 0.583333 & 0.583333\\
\code{astropy__astropy-14508} & Success & Success & 0.785714 & 0.785714\\
\code{pytest-dev__pytest-7432} & Unresolved & Success & 1.000000 & 1.000000\\
\code{pallets__flask-5014} & Success & Success & 1.000000 & 1.000000\\
\code{matplotlib__matplotlib-13989} & Success & Success & 0.500000 & 1.000000\\
\code{psf__requests-2931} & Success & Success & 0.750000 & 0.750000\\
\code{pylint-dev__pylint-6903} & Success & Success & 1.000000 & 1.000000\\
\code{mwaskom__seaborn-3187} & Success & Success & 1.000000 & 0.666667\\
\code{django__django-13513} & Unresolved & Unresolved & 0.647059 & 0.647059\\
\code{sphinx-doc__sphinx-8548} & Unresolved & Unresolved & 0.000000 & 0.000000\\
\code{django__django-13344} & Success & Unresolved & 0.096774 & 0.000000\\
\code{django__django-15098} & Unresolved & Success & 0.000000 & 0.000000\\
\code{django__django-12155} & Unresolved & Unresolved & 0.466667 & 0.466667\\
\code{sphinx-doc__sphinx-8638} & Unresolved & Unresolved & 1.000000 & 1.000000\\
\code{django__django-15252} & Unresolved & Unresolved & 0.000000 & 0.600000\\
\code{django__django-16631} & Success & Success & 0.652174 & 0.652174\\
\code{sphinx-doc__sphinx-8551} & Unresolved & Unresolved & 0.000000 & 0.000000\\
\code{django__django-13810} & Success & Success & 1.000000 & 1.000000\\
\code{sphinx-doc__sphinx-8721} & Unresolved & Unresolved & N/A & N/A\\
\code{sphinx-doc__sphinx-8621} & Unresolved & Unresolved & 0.000000 & 0.000000\\
\code{matplotlib__matplotlib-25311} & Success & Success & 1.000000 & 1.000000\\
\bottomrule
\end{tabular}
\end{table*}


\twocolumn
\section{Reproducibility and Outstanding Information}\label{app:repro}
The host computer is an ASUS Zenbook UX8402VV with a 13th-generation Intel Core i9-13900H processor, 20 logical processors, and approximately 32 GB of RAM. These specifications were verified after the experiment; the saved run records do not independently establish the historical hardware configuration. The workflow used Ubuntu/WSL and Docker, but the CPU and memory allocations available to those environments during execution were not recorded. DeepSeek inference used the hosted API rather than running locally on this hardware.

The evaluation directory is \code{evaluation_results/trace-assigned-26-unit-20261113-122920/}; final selected generation records are in \code{runs/assigned-pairs-20261113-002517/completion-pass-v6/}. Predictions, prompts, patches, trajectories, manifests, image references, workflow hashes, and official pre/post logs are preserved. The dataset revision and harness commit appear in Section~\ref{sec:method}; \code{evidence.json} indexes exact local evidence. Evaluation inputs contain reference fields and must remain outside generation context. No report preparation step modifies either workflow or prediction.

The nominal agent-cost limit excludes preprocessing and does not guarantee a strict cap on spending. Selected preprocessing metrics record 69 requests, 49,789 input tokens, and 173,755 output tokens, including reasoning output reported by the provider. The supplied provider exports, \code{amount-2026-10-18_2026-11-16.csv} and \code{cost-2026-10-18_2026-11-16.csv}, record 1,211 requests for \code{deepseek-flash}, 1,118,057 cache-miss input tokens, 20,423,295 cache-hit input tokens, and 644,720 output tokens. Total billed cost is \$0.959975793. Summing each token quantity multiplied by its listed unit price reproduces both the daily charges and the total exactly. Nonzero charges occur on November 3 and November 10--13. This export-period total may include pilot runs, retries, and other requests; its rows do not identify the method, instance, or individual generation run. We therefore do not allocate it between baseline and TRACE or treat it as the cost of the selected runs alone. No explicit generation seed was configured. The submission date is November 17, 2026. The submission artifacts are archived at commit \code{a5344b0abf54dcb9674f70261ad1a66d938f8a3f}: \url{https://github.com/mkoehle6/trace-swt-bench/tree/a5344b0abf54dcb9674f70261ad1a66d938f8a3f}. The evaluation results are preserved in its \code{evaluation_results/} directory.

\section{Contributions and AI Disclosure}\label{app:disclosure}
Mark Koehler, Hikaru McGuan, and Jason Torres completed this project as a collaborative team effort, sharing responsibility for the research, implementation, evaluation, analysis, and presentation of the work. Author names and email addresses appear on the title page.

The following tools were used in the experiment or supporting project work. Recorded model names are reported as identifiers, not as proof of an unchanged provider backend.

\begin{itemize}
\item \textbf{mini-SWE-agent 2.4.6:} the agent framework that manages the generation conversation, executes shell commands in Docker, and collects submitted test patches. It is not itself a language model. The dependency is pinned in \code{requirements.txt}, and the controller checks this version before launch.
\item \textbf{DeepSeek-Flash:} generated candidate tests and supplied issue extraction and semantic-review responses. The generation configuration records \code{deepseek/deepseek-v4-flash}; the controller's default direct-API analysis identifier is \code{deepseek-flash}. Saved provider responses include the returned model identifier and, where supplied, a system fingerprint. These records do not establish an exact immutable model release for every request.
\item \textbf{OpenAI Codex:} assisted with implementation and workflow orchestration, artifact inspection, evaluation support, analysis, workflow diagrams for presentation slides, and manuscript drafting and revision. The historical Codex application/CLI version and underlying model identifier were not consistently recorded and cannot be specified reliably from the experiment artifacts. A current installation version would not establish the version used for earlier work.
\item \textbf{ChatGPT (GPT-6.1):} supported initial project planning, workflow discussions, troubleshooting, and explanations. The GPT-6.1 model designation was confirmed by the project owner. ChatGPT was not an additional benchmark test-generation arm.
\item \textbf{Claude Sonnet 5.5:} assisted with reviewing the repository, drafting presentation slides, and proofreading the paper; the team reviewed and revised the resulting material. The model designation was confirmed by the project owner.
\end{itemize}

The team confirms that all AI tools used for this project have been disclosed above. This paper does not claim that humans independently validated every generated assertion or verified every action against the experiment's rules. AI-assisted project work and writing are disclosed separately from the evaluated test-generation methods.
